%%=============================================================================
%% Voorwoord
%%=============================================================================

\chapter*{\IfLanguageName{dutch}{Woord vooraf}{Preface}}%
\label{ch:voorwoord}

%% TODO:
%% Het voorwoord is het enige deel van de bachelorproef waar je vanuit je
%% eigen standpunt (``ik-vorm'') mag schrijven. Je kan hier bv. motiveren
%% waarom jij het onderwerp wil bespreken.
%% Vergeet ook niet te bedanken wie je geholpen/gesteund/... heeft

Het schrijven van deze bachelorproef vormt het sluitstuk van mijn opleiding Toegepaste Informatica aan HOGENT. In deze context heb ik gekozen voor het onderwerp automatisatie van het onboardingproces binnen het gemeentebestuur van Berlare. Mijn interesse voor dit onderwerp ontstond vanuit de vaststelling dat veel IT-processen binnen organisaties nog manueel verlopen, terwijl er duidelijke opportuniteiten zijn om deze efficiënter en betrouwbaarder te maken via automatisatie. Deze bevindingen kwamen tot licht tijdens mijn tijd bij de gemeente als student en als stagiair. \newline

Tijdens dit traject kreeg ik de kans om een concreet probleem uit de praktijk te analyseren en om te zetten in een werkbare technische oplossing. Dit maakte het niet alleen een leerrijke, maar ook een bijzonder relevante ervaring met het oog op mijn toekomstige professionele carrière.\newline

Graag wil ik enkele personen bedanken die mij hebben ondersteund bij het tot stand komen van deze bachelorproef. In het bijzonder dank ik mijn promotor, Dhr. Thomas Aelbrecht, voor de begeleiding, feedback en kritische inzichten doorheen het proces. Daarnaast wil ik ook mijn co-promotor, Dhr. Hans Van Mingeroet, bedanken voor de praktische ondersteuning en het delen van zijn expertise binnen het werkveld. Ook wil ik collega Dhr Bram De Wolf bedanken voor zijn praktische ondersteuning. Dankzij Bram en Hans leerde ik de omgeving kennen en kon ik dit project tot een goed einde brengen. \newline

Tot slot wil ik mijn familie en vrienden bedanken voor hun steun en aanmoediging gedurende mijn opleiding en tijdens het schrijven van deze bachelorproef.\newline